\chapter*{Resumen Ejecutivo}
\addcontentsline{toc}{chapter}{Resumen Ejecutivo}

Este informe documenta el proceso ETL aplicado al estudio descriptivo de desigualdad territorial en la Provincia de Santiago, integrando cuatro fuentes públicas sobre áreas verdes comunales, pobreza por ingresos, capacidad municipal e información poblacional. La integración se resolvió sobre la llave \texttt{codigo\_comuna}, debido a que los nombres de comuna difieren entre fuentes en mayúsculas, tildes y estilo de escritura. El repositorio evidencia un flujo reproducible desde perfilado y homologación hasta transformación, integración, validación final, carga a SQLite y análisis exploratorio.

El producto integrado observado en el repositorio corresponde a \texttt{data/processed/desigualdad\_comunal\_final.csv}, con 32 filas y 12 columnas, una fila por comuna y cobertura completa de las 32 comunas del universo definido. La validación formal reporta ausencia de duplicados por \texttt{codigo\_comuna}, ausencia de nulos en las columnas finales, convertibilidad numérica sin errores, rangos consistentes para pobreza, áreas verdes e IPP, y ausencia de infinitos en las métricas derivadas por habitante. La carga final en \texttt{db/lab1\_desigualdad.sqlite} materializa tres tablas: \texttt{dim\_comuna} (32 filas), \texttt{fact\_desigualdad\_comunal} (32 filas) y \texttt{metadata\_fuentes} (4 filas), con \texttt{PRAGMA quick\_check = ok}.

En los outputs de Fase 9 se observan, entre otros resultados descriptivos, menores disponibilidades relativas de áreas verdes en CONCHALI, CERRILLOS, INDEPENDENCIA, LA CISTERNA, SAN MIGUEL; mayores niveles observados de pobreza por ingresos en LA PINTANA, SAN RAMON, LO ESPEJO, EL BOSQUE, CONCHALI; y menores niveles de IPP por habitante en CERRO NAVIA, LA PINTANA, LO PRADO, EL BOSQUE, LA GRANJA. El cruce por quintiles identifica rezago crítico en al menos dos dimensiones para LO ESPEJO, EL BOSQUE, CONCHALI, LA PINTANA, SAN RAMON, LA GRANJA, CERRO NAVIA. Además, Quilicura aparece como outlier plausible en \texttt{areas\_verdes\_m2\_hab} con 3931.614773, muy por encima del siguiente valor observado en La Reina (21.734172), razón por la cual el scatter de pobreza y áreas verdes se documenta con escala logarítmica en el eje X. En conjunto, el dataset final y la base SQLite quedan listos como productos reutilizables para consulta, auditoría metodológica y montaje posterior del informe en el template externo.
